\section{Minimal method revision under observational equivalence}
\label{sec:minimal-revision-theory}

The general object of this paper is not a particular self-editing agent. It is
the decision problem created when several latent failures look the same to a
method reviser but license different scientific changes. The relevant boundary
is therefore informational before it is algorithmic: no more capable proposer,
search procedure or acceptor can recover a distinction that its interface has
erased.

\subsection{Revision fronts and evidence interfaces}

Let $\Theta$ be a set of latent development situations and let
$(\mathcal R,\prec)$ be a well-founded strict partial order of candidate
revisions. The interpretation of $r'\prec r$ is that $r'$ changes strictly less
of the method than $r$; representation repair and evaluator repair may be
incomparable, so a total ladder is neither required nor generally appropriate.
For $\theta\in\Theta$, let $\mathcal A(\theta)\subseteq\mathcal R$ be the
revisions admissible after causal-responsibility, task-recovery, preservation,
fresh-transfer, harm and authority constraints are applied, and let
\[
  \mathcal M(\theta)=\min_{\prec}\mathcal A(\theta)
\]
be its minimal revision front. Well-foundedness guarantees a nonempty minimal
front whenever $\mathcal A(\theta)$ is nonempty.

Let $Q(\theta)$ indicate a material residual. The required decision is the
front-valued map
\[
G(\theta)=
\begin{cases}
\mathsf{NO\_REVISION}, & \neg Q(\theta),\\
\mathsf{UNRESOLVED}, & Q(\theta)\ \text{and}\ \mathcal A(\theta)=\varnothing,\\
\mathcal M(\theta), & Q(\theta)\ \text{and}\ \mathcal A(\theta)\ne\varnothing.
\end{cases}
\]
Returning an antichain rather than silently choosing one incomparable element
is deliberate. Further measurement is required when the evidence does not
select a unique member.

A reviser observes an interface $\phi:\Theta\to\mathcal Z$, not $\theta$
itself. The interface can include a failure transcript, retrieved records,
diagnostic labels, intervention results, issue history, and outputs of any
number of donor systems.

\paragraph{Theorem 1 (revision-factorization theorem).}
There exists an exact revision rule $g:\phi(\Theta)\to\operatorname{im}(G)$
with $G=g\circ\phi$ if and only if
\[
  \phi(\theta)=\phi(\theta')\quad\Longrightarrow\quad
  G(\theta)=G(\theta')
  \qquad\text{for all }\theta,\theta'\in\Theta.
\]

\paragraph{Proof.}
If $G=g\circ\phi$, equal interface values necessarily give equal decisions.
Conversely, if $G$ is constant on each fibre of $\phi$, define $g(z)$ to be
$G(\theta)$ for any $\theta$ satisfying $\phi(\theta)=z$. Fibre constancy
makes the definition independent of the representative. \hfill$\square$

Thus exact selection requires decision sufficiency, not recovery of every
latent cause. A compressed interface is enough if it separates all situations
that require different revision fronts. An arbitrarily rich interface is
insufficient if even one such pair remains observationally equivalent. This is
a decision-specific application of statistical experiment comparison
\cite{blackwell1953}; neither the factorization argument nor Bayes decision
theory below is claimed as a new general statistical principle.

\subsection{Irreducible revision risk}

Let $\Theta\sim\mu$, $Z=\phi(\Theta)$, and $Y=G(\Theta)$, with finite decision
alphabet $\mathcal Y$. Let $\mathcal D$ be a finite action set that may include
revision fronts and abstention, and let
$L:\mathcal D\times\mathcal Y\to[0,\infty)$ encode the cost of under-revision,
false broad revision, harmful transfer, or abstention. Write
$p_y(z)=\Pr\{Y=y\mid Z=z\}$.

\paragraph{Theorem 2 (mixed-fibre Bayes bound).}
Every deterministic or randomized rule using only $Z$ has risk at least
\[
  R^*(\phi,L)=
  \mathbb E_Z\!\left[
    \min_{a\in\mathcal D}\sum_{y\in\mathcal Y}L(a,y)p_y(Z)
  \right],
\]
and a fibrewise Bayes action attains equality. For zero--one loss this becomes
\[
  R^*(\phi)=1-\mathbb E_Z\!\left[\max_y p_y(Z)\right].
\]

\paragraph{Proof.}
Condition on $Z=z$. A randomized rule with probabilities $q(a\mid z)$ has
conditional risk
$\sum_a q(a\mid z)\sum_y L(a,y)p_y(z)$, a convex combination bounded below by
the smallest constituent. Concentrating on a minimizing action attains the
bound. Averaging over $Z$ proves the first claim; zero--one loss yields the
second expression. \hfill$\square$

The bound makes asymmetry explicit. A policy that always chooses the broadest
revision may have low under-revision loss but high false-broad-revision and harm
loss; abstaining may be optimal on an impure fibre without making the fibre
identifiable. No aggregate score can turn missing information into revision
authority.

\paragraph{Corollary 1 (no sound-and-complete self-promotion under authority
equivalence).}
Suppose two situations have the same complete internally visible transcript,
but protected evidence authorizes promotion in one and forbids it in the
other. No promotion rule using only that internal transcript can be both sound
(never promote the forbidden case) and complete (promote the authorized case).

\paragraph{Proof.}
The two situations lie in one interface fibre but require different binary
decisions, contradicting Theorem~1. \hfill$\square$

External custody is therefore not a claim that an external evaluator is
infallible. It is an information-and-authority requirement: the protected
distinction must enter some decision interface not writable by the candidate.
An always-abstain rule can remain sound, but it is not complete.

\subsection{Discriminator design as a covering problem}

Theorems~1--2 say when an interface is insufficient; the next result identifies
exactly what a finite intervention programme must separate. Fix an observed
value $z$ and a finite ambiguous fibre
$F_z=\{\theta:\phi(\theta)=z\}$. Let $\mathcal T$ be candidate tests or
controlled interventions, with deterministic protected outcomes
$o_t(\theta)$ for $t\in\mathcal T$. Define the cross-decision pair universe
\[
 U_z=\bigl\{\{\theta,\theta'\}\subseteq F_z:
 G(\theta)\ne G(\theta')\bigr\}
\]
and the subset distinguished by test $t$,
\[
 D_t=\bigl\{\{\theta,\theta'\}\in U_z:
 o_t(\theta)\ne o_t(\theta')\bigr\}.
\]

\paragraph{Theorem 3 (discriminator-cover theorem).}
A non-adaptive test panel $S\subseteq\mathcal T$ permits exact revision
selection on $F_z$ if and only if
\[
  \bigcup_{t\in S}D_t=U_z.
\]
Consequently, for test costs $c_t\ge0$, a minimum-cost exact panel is exactly
the weighted set-cover optimum of the induced instance
$(U_z,\{D_t:t\in\mathcal T\},c)$. If some pair in $U_z$ belongs to no $D_t$,
no panel drawn from $\mathcal T$ can resolve the revision decision.

\paragraph{Proof.}
The joint signature of $\theta$ under $S$ is
$(o_t(\theta))_{t\in S}$. It supports exact revision selection precisely when
no two situations with different $G$ values share a joint signature. A
cross-decision pair has different joint signatures precisely when at least one
selected test distinguishes it, which is precisely the stated cover condition.
Minimizing the additive test cost over all such panels is the induced weighted
set-cover problem. The last claim follows when a required pair is uncovered by
the full candidate family. \hfill$\square$

\paragraph{Proposition 1 (adaptive leaf-purity criterion).}
For deterministic test outcomes, a terminating adaptive test policy permits an
exact decision on $F_z$ if and only if every reachable terminal leaf contains
only situations with one common $G$ value.

\paragraph{Proof.}
All situations reaching a terminal leaf generate the same observed test path.
An exact terminal decision exists there if and only if their required decisions
agree. Applying this condition to every reachable leaf proves both directions.
\hfill$\square$

These are structural equivalences in the classical minimum-test-set
neighbourhood \cite{moret1985}, not claims that every real experiment is
deterministic or that a generic set-cover heuristic is statistically optimal.
With stochastic outcomes, tests must instead separate protected outcome
distributions to a declared error tolerance. Adaptivity can reduce expected
cost but cannot separate a pair whose outcome distributions are identical under
every available intervention. This connects method
revision to active diagnosis and sequential experimental design
\cite{reiter1987,dekleer1987,chernoff1959,golovin2011} while making the target
the minimal authorized revision, rather than latent-cause recovery alone.

\subsection{The three retained errors as new experiments}

The diagnostic archive contains three observed collisions between neighbouring
revision levels. They are not evidence for Theorems~1--3 and are not repaired
by renaming their labels. Each instead defines a successor pair universe and a
factorial discriminator.

\begin{enumerate}
  \item \textbf{Retrieval versus representation.} Hold the representation
  fixed while varying source access and retrieval policy, then hold retrieved
  information fixed while varying the representation. A retrieval-only rescue,
  a representation-only rescue, an interaction, and a joint null remain four
  distinct outcomes.
  \item \textbf{Environment versus implementation.} Cross a content-addressed
  candidate with independently reconstructed environments. A code defect must
  follow candidate identity across environments; an environmental defect must
  follow environment identity across unchanged code. An interaction remains
  unresolved rather than being forced into either class.
  \item \textbf{Representation versus method basis.} Cross a richer
  representation under a frozen method class with a method-class expansion
  under a frozen representation. Improvement only under the joint treatment
  identifies an interaction, not a license to credit either component alone.
\end{enumerate}

These interventions refine the visible-symptom fibres that produced
P5-HC-002, P5-HC-012 and P5-HC-018. The separate wide successor then asks
whether the resulting decision interface transfers across source-disjoint
project/issue families. A positive local discriminator is therefore a
prerequisite for, not a substitute for, protected fresh-transfer evidence.
